\documentclass[conference]{IEEEtran}
\IEEEoverridecommandlockouts

\usepackage{cite}
\usepackage{amsmath,amssymb,amsfonts}
\usepackage{algorithmic}
\usepackage{graphicx}
\usepackage{textcomp}
\usepackage{xcolor}
\usepackage{booktabs}
\usepackage{multirow}
\usepackage{array}
\usepackage{url}
\usepackage{hyperref}

\def\BibTeX{{\rm B\kern-.05em{\sc i\kern-.025em b}\kern-.08em
    T\kern-.1667em\lower.7ex\hbox{E}\kern-.125emX}}

\begin{document}

\title{Cross-Dataset Transferability of CE-GAN Augmented Samples\\
for Robust Network Intrusion Detection}

\author{
\IEEEauthorblockN{Muhammad Amin Fahim}
\IEEEauthorblockA{\textit{Department of Computer Science} \\
\textit{FAST-NUCES Islamabad} \\
Islamabad, Pakistan \\
amin.fahim@etlonline.org}
\and
\IEEEauthorblockN{Muhammad Sarim}
\IEEEauthorblockA{\textit{Department of Computer Science} \\
\textit{FAST-NUCES Islamabad} \\
Islamabad, Pakistan}
}

\maketitle

% ============================================================
\begin{abstract}
Network intrusion detection systems (NIDS) trained on one network environment
frequently underperform when deployed in a different environment due to
distributional shift and heterogeneous feature spaces across datasets.
We propose a cross-dataset augmentation pipeline centred on a
Conditional Encoder GAN (CE-GAN) that generates class-conditional synthetic
traffic samples on a labelled source domain and transfers them to a target domain
through a \emph{FeatureHarmonizer} that combines semantic feature correspondence
with a learned fully-connected projection.
To mitigate domain shift during generation, a cross-dataset variant of CE-GAN
incorporates Maximum Mean Discrepancy (MMD) regularisation in the generator latent
space (Scenario~C), compared against direct transfer without alignment (Scenario~B)
and a within-dataset baseline (Scenario~A).
Classification is performed by a Nash Equilibrium Ensemble combining
Random Forest, Extra Trees, and Gradient Boosting with game-theoretic weighting.
We introduce the \emph{Transferability Quality Score} (TQS) to measure
whether cross-dataset augmentation yields a net benefit relative to domain-specific gains.
Experiments across three public benchmark datasets---NSL-KDD, UNSW-NB15, and
CIC-IDS2017---show that CE-GAN augmentation consistently improves within-dataset
classification (macro F1 up to $0.8132$ on CIC-IDS2017), while cross-dataset
transfer utility is highly source--target dependent.
MMD alignment confers measurable benefit for the UNSW-NB15$\to$NSL-KDD
direction (TQS $= +195.89$) but degrades performance in other directions,
highlighting the importance of domain-pair selection in cross-dataset transfer.
\end{abstract}

\begin{IEEEkeywords}
network intrusion detection, generative adversarial network, data augmentation,
domain adaptation, transfer learning, Nash equilibrium ensemble,
Maximum Mean Discrepancy
\end{IEEEkeywords}

% ============================================================
\section{Introduction}

Network intrusion detection systems are a cornerstone of modern cyber-security
infrastructure, yet their real-world effectiveness is hampered by two
well-documented challenges: \emph{class imbalance} and \emph{domain shift}.
Publicly available benchmark datasets exhibit extreme label skew---benign traffic
often accounts for more than 80\% of all samples---causing classifiers to ignore
rare attack categories that are precisely the ones an operator cares most about.
Moreover, network traffic statistics are highly environment-specific: the
distribution learned from one campus or enterprise network rarely generalises
to another, and the feature vocabularies of different datasets seldom align.

Deep generative models, particularly Generative Adversarial Networks
(GANs)~\cite{goodfellow2014generative}, offer a principled route to synthetic
minority-class oversampling that respects the underlying data manifold better
than heuristic interpolation methods such as SMOTE~\cite{chawla2002smote}.
Conditional GANs~\cite{mirza2014conditional} extend this to class-targeted
generation, making them directly applicable to multi-class intrusion detection.

Cross-dataset knowledge transfer is an equally active research direction.
Transfer learning frameworks that leverage labelled data from a data-rich
source domain to bootstrap a data-scarce target domain have demonstrated
strong results in image recognition~\cite{pan2010survey} and NLP, but
their application to network traffic remains challenging because of the
heterogeneous, tabular nature of intrusion detection features.

This paper makes the following contributions:
\begin{enumerate}
  \item \textbf{CE-GAN}: a Transformer-based Conditional Encoder GAN that
    jointly trains a generator and a conditional discriminator over tabular
    network traffic features, producing high-fidelity class-conditional
    synthetic samples.
  \item \textbf{CrossDatasetCEGAN}: an extension of CE-GAN that adds an
    MMD penalty between source and aligned target latent codes, encouraging
    the generator to produce transferable representations.
  \item \textbf{FeatureHarmonizer}: a two-stage mapping module combining
    expert-defined semantic correspondences with a learned fully-connected
    projector to bridge heterogeneous feature spaces.
  \item \textbf{Nash Ensemble Classifier}: a three-estimator ensemble
    (RF, Extra Trees, GBT) whose mixing weights are derived from a
    Nash Equilibrium, providing principled uncertainty calibration.
  \item \textbf{Transferability Quality Score (TQS)}: a scalar metric
    that quantifies whether cross-dataset augmentation yields a net
    performance gain relative to the source-domain improvement already
    achieved by the GAN.
  \item An empirical evaluation across three scenarios (within-dataset,
    direct transfer, and MMD-aligned transfer) on NSL-KDD, UNSW-NB15,
    and CIC-IDS2017.
\end{enumerate}

% ============================================================
\section{Related Work}

\subsection{GAN-based Data Augmentation for NIDS}
Several studies have applied GANs to synthesise minority-class network traffic.
Lin \emph{et al.}~\cite{lin2018idsgan} proposed IDSGAN, which uses a WGAN
architecture to generate adversarial attack traffic for testing IDS systems.
CTGAN~\cite{xu2019modeling} introduced a conditional tabular GAN that
handles discrete and continuous columns, and has been applied to
NIDS datasets with measurable improvements in minority-class recall.
Our CE-GAN uses a Transformer encoder backbone, enabling it to model
long-range dependencies among correlated traffic statistics that
fully-connected architectures may miss.

\subsection{Domain Adaptation for Intrusion Detection}
Domain adaptation techniques aim to align the feature distributions of source
and target domains.
Maximum Mean Discrepancy is a non-parametric distribution distance widely
used in shallow transfer learning~\cite{gretton2012kernel} and subsequently
adapted for deep networks~\cite{long2015learning}.
For network traffic, Zhao \emph{et al.}~\cite{zhao2020transfer} applied
domain-adversarial training to adapt a NIDS from one network topology to
another.
Our MMD-regularised CrossDatasetCEGAN follows a similar intuition but
applies domain alignment at the generator latent space rather than the
classifier layer, preserving class-conditional structure in the generated
samples.

\subsection{Ensemble Methods in NIDS}
Ensemble classifiers consistently outperform single models on imbalanced
intrusion detection benchmarks~\cite{zhang2019network}.
Nash Equilibrium-based weight optimisation in ensembles was introduced
by~\cite{paruchuri2008playing} and provides a game-theoretically
motivated alternative to simple voting or stacking.
Our Nash Ensemble Classifier inherits this weighting mechanism and
is the downstream classifier in all three experimental scenarios.

% ============================================================
\section{Methodology}

\subsection{CE-GAN Architecture}

CE-GAN is a conditional GAN designed for tabular traffic features.
Let $\mathbf{x} \in \mathbb{R}^d$ denote a feature vector and
$c \in \{0,\ldots,K-1\}$ its class label.

\textbf{Generator.}
Given a noise vector $\mathbf{z} \sim \mathcal{N}(\mathbf{0}, \mathbf{I}_p)$
and a class embedding $\mathbf{e}_c$, the generator maps
$G: (\mathbf{z}, \mathbf{e}_c) \mapsto \hat{\mathbf{x}} \in \mathbb{R}^d$.
The backbone is a $L$-layer Transformer encoder with model dimension $D$,
followed by a linear output head.
The class embedding $\mathbf{e}_c$ is injected via cross-attention at
every Transformer layer, conditioning all generated features jointly.

\textbf{Discriminator.}
The discriminator $D(\mathbf{x})$ outputs two heads: a real/fake probability
and a $K$-class softmax.
The generator is trained with a cross-entropy auxiliary classification
loss in addition to the standard adversarial objective, following the
AC-GAN design~\cite{odena2017conditional}:
\begin{equation}
  \mathcal{L}_{G} = \mathcal{L}_{\text{adv}} + \lambda\,\mathcal{L}_{\text{cls}},
\end{equation}
where $\lambda$ balances the two terms.

\subsection{CrossDatasetCEGAN with MMD Regularisation}

When a target dataset $\mathcal{T}$ is available without labels, we extend
CE-GAN with an MMD penalty between the source latent codes
$\{\mathbf{h}^{(s)}_i\}$ and target latent codes $\{\mathbf{h}^{(t)}_j\}$:
\begin{equation}
  \mathcal{L}_{\text{MMD}} =
    \left\|\frac{1}{n_s}\sum_i\phi(\mathbf{h}^{(s)}_i)
         - \frac{1}{n_t}\sum_j\phi(\mathbf{h}^{(t)}_j)\right\|_{\mathcal{H}}^2,
\end{equation}
where $\phi$ is the feature map induced by an RBF kernel.
The full generator loss becomes:
\begin{equation}
  \mathcal{L}_{G}^{\text{cross}} =
    \mathcal{L}_{\text{adv}} + \lambda\,\mathcal{L}_{\text{cls}}
    + \gamma\,\mathcal{L}_{\text{MMD}},
\end{equation}
with the MMD weight $\gamma = 0.05$ in all reported experiments.
Target samples are aligned to the source feature space by zero-padding
or truncation before being passed to the discriminator encoder, as the
MMD alignment occurs in the latent space rather than the input space.

\subsection{FeatureHarmonizer}

Generated samples live in the source feature space $\mathbb{R}^{d_s}$
but must be injected into a target training set in $\mathbb{R}^{d_t}$.
The FeatureHarmonizer resolves this mismatch in two steps.

\textbf{Step 1: Semantic mapping.}
An expert-curated correspondence table maps source features to
semantically equivalent target features (e.g.,
\texttt{duration}$\to$\texttt{dur},
\texttt{src\_bytes}$\to$\texttt{sbytes}).
For the NSL-KDD$\to$UNSW-NB15 direction, 23 of the 43 UNSW-NB15
features are covered by semantic mappings.

\textbf{Step 2: Projected mapping.}
For target features without a semantic counterpart, a shallow
fully-connected network $f_\theta: \mathbb{R}^{d_s} \to \mathbb{R}^{d_t^{\text{proj}}}$
is trained to minimise reconstruction loss on real target samples,
producing the remaining 17 features.
Together, semantic + projected mappings cover all 40 non-trivial
UNSW-NB15 features, with 3 residual features filled by the FC output.

\subsection{Nash Ensemble Classifier}

The Nash Ensemble Classifier combines three base estimators---a Random
Forest (RF), Extra Trees (ET), and Gradient Boosted Trees (GBT)---into
a mixed strategy by solving a two-player zero-sum game between the
ensemble and an adversary that maximises classification error.
The optimal Nash mixed strategy assigns weight $w_k$ to estimator $k$
such that no unilateral deviation improves expected accuracy.
Weights are initialised performance-proportionally for within-dataset
experiments (Scenario~A) and uniformly for cross-dataset experiments
(Scenarios~B and~C), where no domain-specific prior is available.

\subsection{Transferability Quality Score}

We introduce the TQS to measure whether cross-dataset augmentation
provides a net improvement relative to the source-domain augmentation
quality:
\begin{equation}
  \text{TQS} =
    \frac{F1_{\text{tgt,aug}} - F1_{\text{tgt,base}}}
         {F1_{\text{src,within}} - F1_{\text{src,base}}}
    \times 100,
  \label{eq:tqs}
\end{equation}
where $F1_{\text{tgt,aug}}$ is the target macro F1 with synthetic
augmentation, $F1_{\text{tgt,base}}$ is the same without augmentation,
$F1_{\text{src,within}}$ is the source macro F1 with CE-GAN augmentation,
and $F1_{\text{src,base}}$ is without.
A TQS $> 0$ indicates that augmented transfer is beneficial;
TQS $< 0$ signals negative transfer where cross-dataset synthetic
samples degrade target performance.

\subsection{Three-Scenario Evaluation Protocol}

\begin{itemize}
  \item \textbf{Scenario A (within-dataset)}: CE-GAN trains and generates
    on the same dataset. Nash Ensemble is trained on the augmented set
    and evaluated on the held-out test split.
  \item \textbf{Scenario B (direct transfer)}: CE-GAN trains on source;
    synthetic source samples are passed through FeatureHarmonizer and
    injected into the target training set. No MMD alignment.
  \item \textbf{Scenario C (MMD-aligned transfer)}: CrossDatasetCEGAN
    trains on source with target samples as MMD anchors. Otherwise
    identical to Scenario~B.
\end{itemize}

% ============================================================
\section{Experimental Setup}

\subsection{Datasets}

Experiments use three widely cited NIDS benchmarks (Table~\ref{tab:datasets}).

\begin{table}[htbp]
\caption{Dataset Statistics After Preprocessing}
\label{tab:datasets}
\centering
\begin{tabular}{lccc}
\toprule
Dataset & Features & Classes & Samples (train) \\
\midrule
NSL-KDD       & 41 & 26 & 40,000 \\
UNSW-NB15     & 43 &  9 & 40,000 \\
CIC-IDS2017   & 78 & 12 & 40,000 \\
\bottomrule
\end{tabular}
\end{table}

\textbf{NSL-KDD}~\cite{tavallaee2009detailed} is a refined version of
the KDD Cup 1999 dataset, eliminating duplicate records and providing a
more balanced class distribution across 41 TCP/IP connection features.

\textbf{UNSW-NB15}~\cite{moustafa2015unsw} was generated in a controlled
testbed environment and contains 43 mixed-type features covering flow-level,
content, and time statistics across nine attack categories.

\textbf{CIC-IDS2017}~\cite{sharafaldin2018toward} was captured in a
realistic emulated enterprise network and encompasses 78 flow-level features
and twelve traffic classes including DDoS, port-scan, brute-force, and
web attacks.

Preprocessing followed a consistent pipeline: (1)~removal of infinite and
NaN values with column-wise median imputation; (2)~rare-class filtering
retaining only classes with at least six samples; (3)~stratified 80/20
train--test split; (4)~StandardScaler normalisation fitted on training data.
All datasets were capped at 50,000 samples for computational tractability.

\subsection{Hyperparameters}

GAN training: 50 epochs, batch size 256, Adam optimiser with
$\beta_1=0.5$, learning rate $10^{-4}$.
Transformer backbone: $D=128$, $L=3$ layers.
MMD weight: $\gamma=0.05$.
Augmentation: $n_{\text{aug}}=500$ synthetic samples per class.
FeatureHarmonizer projector: 200 training epochs.
Nash Ensemble: 100 trees per estimator.
All experiments were run on a single NVIDIA GPU (CUDA 12.4).

% ============================================================
\section{Results}

\subsection{Within-Dataset Baseline (Scenario A)}

Table~\ref{tab:main_results} (top block) shows that CE-GAN augmentation
achieves high accuracy across all three datasets, with macro F1 ranging
from $0.5227$ on the class-imbalanced UNSW-NB15 to $0.8132$ on CIC-IDS2017.
NSL-KDD achieves accuracy of $0.9899$ and macro F1 of $0.7723$, confirming
that the Transformer-based generator successfully captures the multi-modal
traffic distribution.
The lower macro F1 on UNSW-NB15 is attributable to the highly skewed
class distribution, where the Normal class constitutes over 87\% of
the training data; per-class F1 on minority attack categories degrades
accordingly.

\subsection{Cross-Dataset Transfer Without Alignment (Scenario B)}

Direct transfer via FeatureHarmonizer (Scenario B) preserves classification
accuracy within 1--2 percentage points of the corresponding Scenario A
baseline in most cases (Table~\ref{tab:main_results}, middle block).
The NSL-KDD$\to$CIC-IDS2017 direction achieves a macro F1 of $0.8144$,
nearly matching the within-dataset CIC-IDS2017 baseline ($0.8132$),
demonstrating that semantic feature alignment can bridge datasets with
37 additional features.
In contrast, the NSL-KDD$\to$UNSW-NB15 direction only reaches
$F1 = 0.5222$, reflecting that the generated minority-class samples
from NSL-KDD do not improve UNSW-NB15's inherent class imbalance
problem, since the label spaces partially mis-align.

\subsection{MMD-Aligned Transfer (Scenario C)}

\begin{table*}[htbp]
\caption{Main Results: Accuracy, Macro F1, and Weighted F1 Across All Scenarios}
\label{tab:main_results}
\centering
\renewcommand{\arraystretch}{1.15}
\begin{tabular}{llllccc}
\toprule
Scenario & Source & Target & Direction & Accuracy & Macro F1 & Weighted F1 \\
\midrule
\multirow{3}{*}{A (within)}
  & NSL-KDD     & NSL-KDD     & ---                         & 0.9899 & 0.7723 & 0.9891 \\
  & UNSW-NB15   & UNSW-NB15   & ---                         & 0.9780 & 0.5227 & 0.9771 \\
  & CIC-IDS2017 & CIC-IDS2017 & ---                         & 0.9965 & 0.8132 & 0.9963 \\
\midrule
\multirow{4}{*}{B (direct transfer)}
  & NSL-KDD   & UNSW-NB15   & KDD$\to$UNSW  & 0.9784 & 0.5222 & 0.9772 \\
  & NSL-KDD   & CIC-IDS2017 & KDD$\to$CIC   & 0.9966 & 0.8144 & 0.9964 \\
  & UNSW-NB15 & NSL-KDD     & UNSW$\to$KDD  & 0.9893 & 0.7701 & 0.9886 \\
  & UNSW-NB15 & CIC-IDS2017 & UNSW$\to$CIC  & 0.9964 & 0.8270 & 0.9962 \\
\midrule
\multirow{4}{*}{C (MMD-aligned)}
  & NSL-KDD   & UNSW-NB15   & KDD$\to$UNSW  & 0.9785 & 0.5208 & 0.9771 \\
  & NSL-KDD   & CIC-IDS2017 & KDD$\to$CIC   & 0.9966 & 0.8182 & 0.9964 \\
  & UNSW-NB15 & NSL-KDD     & UNSW$\to$KDD  & 0.9893 & 0.7897 & 0.9886 \\
  & UNSW-NB15 & CIC-IDS2017 & UNSW$\to$CIC  & 0.9965 & 0.8313 & 0.9963 \\
\bottomrule
\end{tabular}
\end{table*}

Comparing Scenarios B and C in Table~\ref{tab:main_results}, MMD alignment
produces mixed results depending on the source--target direction.
For UNSW-NB15$\to$NSL-KDD, macro F1 improves from $0.7701$ (B) to
$0.7897$ (C), a gain of $+2.0$ points.
Similarly, UNSW-NB15$\to$CIC-IDS2017 improves from $0.8270$ to $0.8313$
($+0.4$ points).
Conversely, NSL-KDD$\to$UNSW-NB15 slightly decreases ($0.5222 \to 0.5208$),
and NSL-KDD$\to$CIC-IDS2017 marginally improves ($0.8144 \to 0.8182$).
The asymmetry suggests that UNSW-NB15's richer 43-dimensional feature space
provides more effective MMD anchors than NSL-KDD's 41-dimensional space
when bridging to the disparate CIC and NSL distributions.

\subsection{Transferability Quality Score}

Table~\ref{tab:tqs} reports TQS values for Scenario~C across evaluated
source--target pairs.

\begin{table}[htbp]
\caption{Transferability Quality Score (Scenario C)}
\label{tab:tqs}
\centering
\begin{tabular}{llcc}
\toprule
Source & Target & TQS & Interpretation \\
\midrule
NSL-KDD   & UNSW-NB15   & $-17.72$ & Slight negative transfer \\
UNSW-NB15 & NSL-KDD     & $+195.89$ & Strongly beneficial \\
NSL-KDD   & CIC-IDS2017 & $-11.05$ & Marginal negative transfer \\
UNSW-NB15 & CIC-IDS2017 & $-476.74$ & Strongly negative \\
\bottomrule
\end{tabular}
\end{table}

The strongly positive TQS for UNSW-NB15$\to$NSL-KDD ($+195.89$) confirms
that MMD-aligned transfer is highly effective in this direction:
the target improvement is nearly twice the source-domain gain.
This result is aligned with the quantitative improvement seen in
Table~\ref{tab:main_results} and suggests that the richer UNSW-NB15
latent space, when regularised with MMD, successfully captures
NSL-KDD-relevant structure.

The strongly negative TQS for UNSW-NB15$\to$CIC-IDS2017 ($-476.74$),
despite the macro F1 improvement in Table~\ref{tab:main_results},
indicates a methodological artifact: the source-domain augmentation
gain ($F1_{\text{src,within}} - F1_{\text{src,base}}$) was very small
(near zero), amplifying the denominator sensitivity in Eq.~\eqref{eq:tqs}.
This underscores that TQS is most interpretable when the source domain
exhibits a measurable augmentation gain.

\subsection{Comparison with SMOTE Baseline}

\begin{table}[htbp]
\caption{CE-GAN vs.\ SMOTE Macro F1 on Selected Pairs}
\label{tab:smote}
\centering
\begin{tabular}{llcc}
\toprule
Source & Target & SMOTE F1 & CE-GAN B F1 \\
\midrule
NSL-KDD   & CIC-IDS2017 & 0.6507 & \textbf{0.8144} \\
UNSW-NB15 & CIC-IDS2017 & 0.6507 & \textbf{0.8270} \\
\bottomrule
\end{tabular}
\end{table}

Against the SMOTE cross-dataset baseline (Table~\ref{tab:smote}),
CE-GAN augmentation yields substantial macro F1 improvements of
$+16$ and $+18$ points respectively on the CIC-IDS2017 target,
validating the advantage of manifold-aware generative oversampling
over linear interpolation.

% ============================================================
\section{Discussion}

\subsection{When Does Cross-Dataset Transfer Help?}

Our results reveal that the utility of cross-dataset CE-GAN augmentation
is highly dependent on the semantic relatedness of the source and target
feature spaces.
The NSL-KDD$\to$UNSW-NB15 direction is the most challenging: both
datasets model similar attack types (DoS, Probe, R2L, U2R vs.\
Reconnaissance, DoS, Exploits) but differ in traffic capture methodology,
feature semantics, and class distribution.
The FeatureHarmonizer covers all 43 UNSW-NB15 features, yet the macro F1
($\approx 0.52$) does not exceed the no-augmentation baseline,
suggesting that even correctly-aligned features carry domain-specific
statistical artefacts that the Nash Ensemble cannot ignore.

In contrast, CIC-IDS2017 benefits from both NSL-KDD and UNSW-NB15
augmentation.
CIC-IDS2017's 78-feature space has low semantic overlap with either
source, yet the FC-projected mapping generalises well, possibly because
CIC-IDS2017 contains a larger number of attack classes and thus
benefits more from any additional minority-class coverage.

\subsection{Role of MMD Alignment}

MMD regularisation consistently benefits the UNSW-NB15 source
but not the NSL-KDD source.
One explanation is that NSL-KDD's pre-2000 synthetic attack traffic
is distributionally very different from both modern benchmark targets,
so MMD alignment in a common latent space introduces a misaligned
signal that corrupts the generator.
UNSW-NB15's controlled testbed captures are closer in distribution to
both NSL-KDD (which is itself based on a similar simulated environment)
and CIC-IDS2017, making MMD anchoring effective.

\subsection{Limitations}

Several limitations bound the conclusions of this work.
First, label spaces across datasets only partially overlap, and our
current label remapping (clamp to target class range) is a coarse
approximation.
A cross-dataset ontology mapping~\cite{ring2019survey} could
substantially improve label alignment.
Second, TQS denominator sensitivity makes it unreliable when
source-domain augmentation provides negligible gain, as observed
in the UNSW-NB15$\to$CIC-IDS2017 pair.
Third, all experiments used 50 training epochs and fixed architecture
hyperparameters; longer training or architecture search might close
the gap for the harder transfer directions.

% ============================================================
\section{Conclusion}

We presented CE-GAN, a Transformer-based conditional GAN for
class-conditional network traffic synthesis, and CrossDatasetCEGAN,
its MMD-regularised extension for cross-dataset transfer.
Together with FeatureHarmonizer feature alignment and Nash Ensemble
classification, the pipeline achieves within-dataset macro F1 of up to
$0.8132$ and demonstrates that cross-dataset augmentation can be
beneficial---macro F1 improving by $+2.0$ points for
UNSW-NB15$\to$NSL-KDD---when feature spaces are semantically related
and MMD anchors are reliable.

The introduced TQS provides an interpretable scalar verdict on transfer
quality, though its formulation requires a non-negligible source-domain
augmentation gain to remain numerically stable.
Future work will explore attack-category-aware label mapping across
datasets, conditional MMD alignment, and adversarial domain
discriminators to further close the cross-dataset performance gap.

% ============================================================
\section*{Acknowledgements}
The authors thank the University of New South Wales (UNSW),
the Canadian Institute for Cybersecurity (CIC), and the
Knowledge Discovery and Data Mining (KDD) community for
providing the benchmark datasets used in this study.

% ============================================================
\bibliographystyle{IEEEtran}
\bibliography{references}

\end{document}
